% Requires: tcolorbox, tikz, xcolor (with darkred, darkgreen defined in preamble)

\begin{figure}[ht]
\centering
\begin{tikzpicture}[
    box/.style={draw, rounded corners=3pt, minimum width=3.8cm, minimum height=1.2cm, align=center, font=\small},
    arrow/.style={->, thick, >=stealth},
    label/.style={font=\footnotesize\itshape, text width=2.5cm, align=center}
]

% ============ TRADITIONAL PROGRAMMING (Left) ============
\node[font=\bfseries\sffamily, text=darkred] at (-4.5, 3.2) {Traditional Programming};

% Input boxes
\node[box, fill=gray!10] (data1) at (-5.8, 1.8) {Data};
\node[box, fill=red!10, draw=darkred] (rules1) at (-3.2, 1.8) {Rules\\{\footnotesize (written by you)}};

% Computer box
\node[box, fill=blue!8, draw=blue!40!black, minimum width=4cm] (computer1) at (-4.5, 0) {Computer};

% Output box
\node[box, fill=green!10, draw=darkgreen] (output1) at (-4.5, -1.8) {Answers};

% Arrows
\draw[arrow] (data1.south) -- ++(0,-0.4) -| (computer1.north west);
\draw[arrow] (rules1.south) -- ++(0,-0.4) -| (computer1.north east);
\draw[arrow] (computer1.south) -- (output1.north);

% Labels
\node[label, anchor=east] at (-6.8, 1.8) {input};
\node[label, anchor=west] at (-2.2, 1.8) {input};
\node[label, anchor=east] at (-6.2, -1.8) {output};

% ============ DIVIDER ============
\draw[thick, dashed, gray] (0, 3.5) -- (0, -2.8);
\node[font=\Large\bfseries, gray] at (0, -3.2) {vs};

% ============ MACHINE LEARNING (Right) ============
\node[font=\bfseries\sffamily, text=darkgreen] at (4.5, 3.2) {Machine Learning};

% Input boxes
\node[box, fill=gray!10] (data2) at (3.2, 1.8) {Data};
\node[box, fill=green!10, draw=darkgreen] (answers2) at (5.8, 1.8) {Answers\\{\footnotesize (examples)}};

% Computer box
\node[box, fill=blue!8, draw=blue!40!black, minimum width=4cm] (computer2) at (4.5, 0) {Computer};

% Output box
\node[box, fill=red!10, draw=darkred] (rules2) at (4.5, -1.8) {Rules\\{\footnotesize (learned model)}};

% Arrows
\draw[arrow] (data2.south) -- ++(0,-0.4) -| (computer2.north west);
\draw[arrow] (answers2.south) -- ++(0,-0.4) -| (computer2.north east);
\draw[arrow] (computer2.south) -- (rules2.north);

% Labels
\node[label, anchor=east] at (2.2, 1.8) {input};
\node[label, anchor=west] at (6.8, 1.8) {input};
\node[label, anchor=west] at (6.2, -1.8) {output};

\end{tikzpicture}

\vspace{1em}

\begin{tcolorbox}[
    colback=blue!3!white,
    colframe=blue!40!black,
    arc=2pt,
    boxrule=0.5pt,
    width=0.95\textwidth,
    halign=center,
    top=3pt, bottom=3pt
]
\textbf{Key insight:} In traditional programming, \emph{you} write the rules. In machine learning, the algorithm \emph{discovers} the rules by studying examples.
\end{tcolorbox}

\caption{Traditional programming versus machine learning. On the left, the programmer supplies both data and hand-coded rules; the computer produces answers. On the right, the programmer supplies data and example answers; the computer produces the rules (a trained model).}
\label{fig:traditional_vs_ml}
\end{figure}